% Traduction française COMIMa de l'énoncé original (anglais).
On considère les cases unités du quadrillage $(x, y)$, avec $x, y$ entiers positifs ou nuls, chacune repérée par son coin inférieur gauche. On grise les cases $(0, 0)$, $(1, 0)$, $(0, 1)$, $(2, 0)$, $(1, 1)$, $(0, 2)$. (a) Il y a un jeton sur chacune de ces six cases ; (b) il n'y a qu'un jeton, sur $(0, 0)$.

\emph{Étape} : si $(x, y)$ est occupée et que $(x + 1, y)$ et $(x, y + 1)$ sont libres, on peut retirer le jeton de $(x, y)$ et poser un jeton sur chacune des cases $(x + 1, y)$ et $(x, y + 1)$. Le but est de libérer toutes les cases grisées. Est-ce possible dans le cas (a) ? dans le cas (b) ? (Kontsevich, TT 1981.)
